% Back matter: rules questions, glossary, sources
\dosdivider{07}{Questions, terms and sources}{The rules questions worth asking in round one, the vocabulary of this dossier, and every source it rests on.}
  {\dosdivrow{0}{10}{Rules Q\&A, ranked by points at stake}\dosdivrow{1}{11}{Glossary}\dosdivrow{2}{12}{Sources}}

\section{Rules Q\&A, ranked by points at stake}\label{sec:qa}
\subsection{How to use this list}
Round one of the official Q\&A opens on 1 November and accepts five questions; later rounds accept three. Each question below states what rides on the answer and the design the team builds to until it is answered, so no answer can strand the aircraft.

\begin{widetable}[h]{Round-one questions, in order of the points they protect. The fallback is what the design assumes until the answer arrives.}{tab:qa}
\begin{brieftable}{width=\fullw, colspec={Q[l,wd=0.35in] X[1.45,l] X[1.0,l] X[1.1,l]}}
Rank & Question, as it would be submitted & What rides on it & Fallback until answered \\
1 & In Mission~3, is the time to stow the sensor after the final lap inside the 300~s flight window? & Seven laps against six at 39~s laps; the choice of drum drive & Assume it counts; trim the Mission~3 ballast so the stow fits \\
2 & Staging allows no adjustments such as opening hatches. May a hatch be opened and latched inside the box, without tools, to load payload and batteries? & A misreading fails staging & Enter with the bay open; latch-only closure inside the box \\
3 & Does the 381~mm ring-to-bottom limit apply to a sling to the container's corners, and may the sling legs sit at 10 to 15°? & The container envelope and the declared drop height & 150~g sensor and 41~mm foam fit at 15°; a lower drop height costs 8\,\% of the Ground Mission bonus per 5~in \\
4 & Do identical bolt-on spare panels and control surfaces count as repair of the original airframe? & Whether one broken surface ends the contest & Carry identical bolt-on spares as repair stock \\
5 & May the sensor's internal battery stay in the sensor while it is packed in its container for the Ground Mission? & Ground Mission compliance & A removable sensor battery fitted at the light check \\
Later & Is the 1.5-span deployment distance checked statically or in flight? & Line length & Static exit-to-nose length with a 0.10~m margin \\
Later & Is wingspan the projected tip-to-tip width when vertical tip surfaces are fitted? & Only matters if winglets return & Not needed: winglets were eliminated (Table~\ref{tab:elim}) \\
\end{brieftable}
\end{widetable}

\railnote{Already answered}{Packs may change between flights, within the identical-pack rule (FAQ, layout lines 1670--1673). Rate gyros and stabilisation are allowed (line 600). Neither needs a question.}

\section{Glossary}
\begin{railtable}[h]{Terms used in this dossier.}{tab:gloss}
\begin{brieftable}{width=\textwidth, colspec={Q[l,wd=1.2in] X[l]}}
Term & Meaning \\
Declared sensor mass & The maximum sensor mass the team states at technical inspection; it fixes every Mission~2 unit and the Ground Mission score \\
Payload unit & One container holding the sensor at its declared mass, or one simulator matching it within 28~g \\
Field maximum & The best result any team posts in a mission, our own included; every bonus is a ratio to it \\
P(top) & Probability that our total mission score is the highest in the field \\
P(1st) & Probability that our competition score (report times mission, plus participation) is the highest \\
P10, P90 & The score below which the worst tenth, or above which the best tenth, of the uncertain outcomes fall \\
Paired draws & The same 120 sets of uncertain inputs applied to every candidate, so differences are real \\
Trimmed $C_{L,\max}$ & Maximum lift after the tail's balancing load is subtracted \\
Oswald $e$ & How close the wing's drag due to lift is to the ideal; 1.0 is ideal \\
Fowler flap & A flap that slides aft on tracks as it lowers, adding wing area and opening a slot \\
Single-slotted flap & A flap that rotates down and opens a gap that re-energises the flow over it \\
Stall speed $V_s$ & The slowest speed at which the wing can carry the aircraft's weight; landings are made somewhat above it \\
Decision timing & Choosing each number at the moment the information to choose it well exists, and not before \\
\end{brieftable}
\end{railtable}

\section{Sources}
\subsection{How sources are graded}
Each row gives the source, what was taken from it, when it was read, and its nature: measured data, a specification, a computation, a rule, or a method. Sources used by DB-01 and the analysis scripts are indexed in the repository's \texttt{SOURCES.md}.

% generated by report/make_sources.py; do not edit
\begin{dbflongtable}{colspec={Q[l,wd=0.95in] X[3,l] Q[l,wd=0.7in]}, width=\textwidth}
Source & Reference and use & Nature \\
{\small\ttfamily REF} & AIAA (2026). \emph{2026-\allowbreak{}27 Design/\allowbreak{}Build/\allowbreak{}F\allowbreak{}ly Competition Rules (draft of 2 September 2026)}. {\small\color{dbfSteel}rules/\allowbreak{}2027\_DBF\allowbreak{}\_Rules\_Draft.p\allowbreak{}df; SHA-\allowbreak{}256 d94773fb...ae1}\newline{\small All rule quotes, by line of the layout text} & {\small PRIMARY} \\
{\small\ttfamily REF} & AIAA (2026). \emph{DBF competition summaries and top-\allowbreak{}3 design reports, 2021-\allowbreak{}26}. {\small\color{dbfSteel}aiaa.org; registry analysis/\allowbreak{}data/\allowbreak{}\allowbreak{}sources.toml (1,008 verified values)}\newline{\small Field data, mass, drag, lap times} & {\small PRIMARY} \\
{\small\ttfamily REF} & Team (DB-\allowbreak{}01) (2026). \emph{DB-\allowbreak{}01 Sizing and configuration decision brief, revision A, 14 October 2026}. {\small\color{dbfSteel}decisions/\allowbreak{}DB-\allowbreak{}0\allowbreak{}1\_sizing\_and\_c\allowbreak{}onfiguration.p\allowbreak{}df}\newline{\small Baseline and control design} & {\small PRIMARY} \\
{\small\ttfamily REF} & Team (analysis) (2026). \emph{Analysis Scripts 1-\allowbreak{}6 and 3b; SOURCES.md registry (113 sources)}. {\small\color{dbfSteel}analysis/\allowbreak{}; SOURCES.md}\newline{\small Scoring, sizing, propulsion and Monte Carlo models} & {\small PRIMARY} \\
{\small\ttfamily S1-\allowbreak{}RULES-\allowbreak{}2026-\allowbreak{}\allowbreak{}09-\allowbreak{}02} & AIAA (2026). \emph{2026-\allowbreak{}27 DBF rules, DRAFT dated 2 Sep 2026}. {\small\color{dbfSteel}REPO/\allowbreak{}rules/\allowbreak{}202\allowbreak{}7\_DBF\_Rules\_Dr\allowbreak{}aft.pdf, SHA-\allowbreak{}256 d94773fb17ee3b\allowbreak{}dac88191bf4cbc\allowbreak{}e96c8d0b2166b9\allowbreak{}b089ed670052bf\allowbreak{}9ca3bae1; read 2026-\allowbreak{}10-\allowbreak{}27 and 2026-\allowbreak{}10-\allowbreak{}29}\newline{\small all 99 clause quotes in out/\allowbreak{}rules/\allowbreak{}comp\allowbreak{}liance.csv, checked against analysis/\allowbreak{}data/\allowbreak{}\allowbreak{}raw/\allowbreak{}aiaa-\allowbreak{}dbf-\allowbreak{}2\allowbreak{}027-\allowbreak{}rules-\allowbreak{}draf\allowbreak{}t.layout.txt by tests/\allowbreak{}test\_v\_r\allowbreak{}ules.py} & {\small RULE} \\
{\small\ttfamily S1-\allowbreak{}RESULTS-\allowbreak{}202\allowbreak{}2-\allowbreak{}26} & AIAA, parsed in REPO/\allowbreak{}analysis/\allowbreak{}\allowbreak{}data/\allowbreak{}dbf\_resul\allowbreak{}ts.csv (2026). \emph{AIAA DBF competition results 2022-\allowbreak{}2026 (513 team rows)}. {\small\color{dbfSteel}REPO/\allowbreak{}analysis/\allowbreak{}\allowbreak{}data/\allowbreak{}dbf\_resul\allowbreak{}ts.csv; read 2026-\allowbreak{}10-\allowbreak{}29}\newline{\small report and mission score distributions, winning scores 501.26-\allowbreak{}567.51, participation 3} & {\small MEASURED} \\
{\small\ttfamily S1-\allowbreak{}ATTEMPTS} & REPO/\allowbreak{}analysis/\allowbreak{}\allowbreak{}data/\allowbreak{}attempts.\allowbreak{}py (AIAA competition summaries 2022-\allowbreak{}26) (2026). \emph{Attempt-\allowbreak{}budget model, 2.4-\allowbreak{}2.9 official attempts per team that passed tech, 57-\allowbreak{}73 \% of attempts scored}. {\small\color{dbfSteel}REPO/\allowbreak{}analysis/\allowbreak{}\allowbreak{}data/\allowbreak{}attempts.\allowbreak{}py; read 2026-\allowbreak{}10-\allowbreak{}29}\newline{\small p\_complete for the M1 waiver value} & {\small COMPUTED} \\
{\small\ttfamily S1-\allowbreak{}ADAPTER} & S12 (this effort) (2026). \emph{S12 integration adapter over DB-\allowbreak{}01 s6\_studies.eva\allowbreak{}luate}. {\small\color{dbfSteel}W/\allowbreak{}src/\allowbreak{}integrat\allowbreak{}ion/\allowbreak{}adapter.py; read 2026-\allowbreak{}10-\allowbreak{}29}\newline{\small H1 study: 120 paired draws seed 2027, expected field; REF-\allowbreak{}D01 6.3326} & {\small COMPUTED} \\
{\small\ttfamily S3-\allowbreak{}naca-\allowbreak{}tr534} & R. C. Platt (1935). \emph{Aerodynamic characteristic\allowbreak{}s of a wing with Fowler flaps, including flap loads, downwash, and calculated effect on take-\allowbreak{}off}. {\small\color{dbfSteel}NACA Report 534 (NTRS); local copy out/\allowbreak{}aero\_mono/\allowbreak{}\allowbreak{}raw/\allowbreak{}naca/\allowbreak{}tr534\allowbreak{}.pdf; read 2026-\allowbreak{}10-\allowbreak{}28}\newline{\small fig. 5 (p. 4): CLmax of a Clark Y wing AR 6, Re 6.1e5, with 0.20c/\allowbreak{}0.30c/\allowbreak{}0.\allowbreak{}40c full-\allowbreak{}span Fowler flaps, 1.29 clean to 2.45/\allowbreak{}2.83/\allowbreak{}3.17 on retracted area; Tables II-\allowbreak{}III (p. 15): Cm c/\allowbreak{}4 -\allowbreak{}0.085 clean, -\allowbreak{}0.508 at alpha 0 with 0.20c flap at 30 deg} & {\small MEASURED} \\
{\small\ttfamily S3-\allowbreak{}naca-\allowbreak{}tr664} & C. J. Wenzinger, T. A. Harris (1939). \emph{Wind-\allowbreak{}tunnel investigation of an N.A.C.A. 23012 airfoil with various arrangements of slotted flaps}. {\small\color{dbfSteel}NACA Report 664 (NTRS); local copy out/\allowbreak{}aero\_mono/\allowbreak{}\allowbreak{}raw/\allowbreak{}naca/\allowbreak{}tr664\allowbreak{}.pdf; read 2026-\allowbreak{}10-\allowbreak{}28}\newline{\small fig. 30: increments of section Clmax for split, plain, external-\allowbreak{}airfo\allowbreak{}il, Fowler and slotted 2-\allowbreak{}h flaps vs deflection (read by eye, +/\allowbreak{}-\allowbreak{}0.05); p. 29: slotted-\allowbreak{}flap pitching moment equals plain-\allowbreak{}flap moment at equal lift} & {\small MEASURED} \\
{\small\ttfamily S3-\allowbreak{}leng2022} & Leng et al. (2022). \emph{A study of high-\allowbreak{}lift devices for low Reynolds number applications}. {\small\color{dbfSteel}ICAS 2022; local copy out/\allowbreak{}aero\_mono/\allowbreak{}\allowbreak{}raw/\allowbreak{}icas2022/\allowbreak{}l\allowbreak{}eng2022.pdf; read 2026-\allowbreak{}10-\allowbreak{}28}\newline{\small Sec. 6.3-\allowbreak{}6.4: DAE31, Re 2.3e5, clean CLmax 1.48; single-\allowbreak{}slotted 30 \% flap at 30 deg CLmax 2.5 (+1.0); ESDU estimate +0.587; DATCOM 1.10} & {\small MEASURED (quasi-\allowbreak{}2-\allowbreak{}D tunnel with end plates)} \\
{\small\ttfamily S3-\allowbreak{}williamson2\allowbreak{}012} & G. Williamson (2012). \emph{Plain flaps on low-\allowbreak{}Re airfoils (UIUC MS thesis); UIUC LSAT Vol 6}. {\small\color{dbfSteel}REPO analysis/\allowbreak{}data/\allowbreak{}\allowbreak{}raw/\allowbreak{}aero/\allowbreak{}willi\allowbreak{}amson-\allowbreak{}2012-\allowbreak{}the\allowbreak{}sis.pdf and uiuc-\allowbreak{}lsat-\allowbreak{}v6-\allowbreak{}d\allowbreak{}ata.zip; read 2026-\allowbreak{}10-\allowbreak{}29}\newline{\small AG40d-\allowbreak{}02r 25 \% plain flap: dCl, dCm, dClmax vs deflection 5-\allowbreak{}40 deg at Re 3e5; flap drag files} & {\small MEASURED} \\
{\small\ttfamily S3-\allowbreak{}uiuc-\allowbreak{}lsat-\allowbreak{}s\allowbreak{}d7037} & Selig et al. (1995). \emph{UIUC Low-\allowbreak{}Speed Airfoil Tests Vol 1-\allowbreak{}3, SD7037 drag files (DRAG01-\allowbreak{}03.TXT\allowbreak{})}. {\small\color{dbfSteel}REPO analysis/\allowbreak{}data/\allowbreak{}\allowbreak{}raw/\allowbreak{}aero/\allowbreak{}uiuc-\allowbreak{}\allowbreak{}lsat-\allowbreak{}v1-\allowbreak{}data.z\allowbreak{}ip (also v2, v3); read 2026-\allowbreak{}10-\allowbreak{}29}\newline{\small SD7037 cd at Re about 3e5: clean cdmin 0.0071-\allowbreak{}0.0077; 0.0092-\allowbreak{}0.0097 at cl 0.7; 0.0131-\allowbreak{}0.0145 at cl 1.0; tripped 0.0125 /\allowbreak{} 0.0178 /\allowbreak{} 0.0233} & {\small MEASURED} \\
{\small\ttfamily S3-\allowbreak{}xfoil} & computed with REPO analysis/\allowbreak{}tools\allowbreak{}/\allowbreak{}xfoil.exe via analysis/\allowbreak{}xfoil\allowbreak{}\_tools.py (import only) (2026). \emph{XFoil 6.99 polars: SD7037 Re 3e5-\allowbreak{}1.1e6 Ncrit 9 and 5; SD7034, E387, MH114 Re 5e5 and 8e5}. {\small\color{dbfSteel}; read 2026-\allowbreak{}10-\allowbreak{}29}\newline{\small profile drag a + kp CL\^{}2; section comparison; XFoil CLmax over-\allowbreak{}predicts (SD7037 Re 3e5: 1.38 vs 1.22-\allowbreak{}1.30 measured)} & {\small COMPUTED} \\
{\small\ttfamily S3-\allowbreak{}avl} & M. Drela, H. Youngren (2025). \emph{AVL 3.52 vortex lattice: Trefftz-\allowbreak{}plane e, critical-\allowbreak{}secti\allowbreak{}on loading, winglet, flap-\allowbreak{}twist and biplane studies}. {\small\color{dbfSteel}; read 2026-\allowbreak{}10-\allowbreak{}29}\newline{\small e of tapered wings 0.988-\allowbreak{}0.999; winglet de; biplane k ratio 0.77 at gap/\allowbreak{}span 0.15 (literature about 0.78, Prandtl 1924, not re-\allowbreak{}read)} & {\small COMPUTED} \\
{\small\ttfamily S3-\allowbreak{}kroo2000} & I. Kroo (2001). \emph{Drag due to lift: concepts for prediction and reduction}. {\small\color{dbfSteel}Annu. Rev. Fluid Mech. 33, 587-\allowbreak{}617; read not re-\allowbreak{}read this session; method relayed by the orchestrator 29 Oct 2026}\newline{\small fuselage factor 1 -\allowbreak{} 2 (d/\allowbreak{}b)\^{}2; K\_visc about 0.38 CD0} & {\small METHOD} \\
{\small\ttfamily S3-\allowbreak{}hoerner} & S. F. Hoerner (1965). \emph{Fluid-\allowbreak{}Dynamic Drag}. {\small\color{dbfSteel}; read not re-\allowbreak{}read; standard formulas}\newline{\small fuselage form factor 1 + 1.5 (d/\allowbreak{}l)\^{}1.5 + 7 (d/\allowbreak{}l)\^{}3} & {\small METHOD} \\
{\small\ttfamily S3-\allowbreak{}tn663-\allowbreak{}cal} & REPO analysis/\allowbreak{}wing\_\allowbreak{}clmax\_validati\allowbreak{}on.py (2026). \emph{NACA TN 663 area-\allowbreak{}method calibration, kappa 1.086 (clean) and 1.075 (flap increment)}. {\small\color{dbfSteel}; read 2026-\allowbreak{}10-\allowbreak{}29}\newline{\small area method calibration and range} & {\small COMPUTED} \\
{\small\ttfamily S3-\allowbreak{}gear} & REPO analysis/\allowbreak{}out/\allowbreak{}c\allowbreak{}onfig\_trades\_r\allowbreak{}eport.txt; claude\_docs/\allowbreak{}ev\allowbreak{}idence\_plan.md (E06) (2026). \emph{GT 2026 measured fixed-\allowbreak{}gear drag area 9.9 cm2, Joanneum 2022 predicted 19.4 cm2; NACA TR 485 fairing cuts drag about 23 \%}. {\small\color{dbfSteel}; read 2026-\allowbreak{}10-\allowbreak{}29}\newline{\small gear drag area range 9.9-\allowbreak{}19.4 cm2} & {\small MEASURED /\allowbreak{} COMPUTED} \\
{\small\ttfamily S3-\allowbreak{}rq4-\allowbreak{}cfd} & S10 /\allowbreak{} REPO cfd/\allowbreak{}cases/\allowbreak{}rq4\_\allowbreak{}mock (2026). \emph{rq4\_mock CFD RESULTS.md (indicative only)}. {\small\color{dbfSteel}; read 2026-\allowbreak{}10-\allowbreak{}29}\newline{\small non-\allowbreak{}wing drag area about 0.0070 m2 at alpha 0 as a cross-\allowbreak{}check, never a calibration} & {\small COMPUTED} \\
{\small\ttfamily S6-\allowbreak{}LF-\allowbreak{}MAXI32} & Littelfuse (2026). \emph{MAXI Blade Fuse Rated 32V, series 0299 (datasheet)}. {\small\color{dbfSteel}https:/\allowbreak{}/\allowbreak{}static\allowbreak{}.chipdip.ru/\allowbreak{}li\allowbreak{}b/\allowbreak{}649/\allowbreak{}DOC02864\allowbreak{}9404.pdf (mirror; littelfuse.com refused scripted access); read 2026-\allowbreak{}10-\allowbreak{}29}\newline{\small voltage rating 32 V DC, 1000 A interrupting rating at 32 V DC, ratings 20-\allowbreak{}80 A} & {\small SPEC} \\
{\small\ttfamily S6-\allowbreak{}LF-\allowbreak{}SMARTGLO\allowbreak{}W} & Littelfuse; Digi-\allowbreak{}Key (2023). \emph{Indicator MAXI Smart Glow fuse, Rated 32 V DC, rev 5/\allowbreak{}03/\allowbreak{}2023; Digikey listings 0MAX100.HXGLO and 0MAX100.XPGLO}. {\small\color{dbfSteel}https:/\allowbreak{}/\allowbreak{}www.di\allowbreak{}gikey.com (keyword 0MAX100); datasheet via Digikey-\allowbreak{}hosted copy; read 2026-\allowbreak{}10-\allowbreak{}29}\newline{\small the only 100 A MAXI blade found is rated 32 V DC (value F-\allowbreak{}FUSE-\allowbreak{}100-\allowbreak{}32)} & {\small SPEC} \\
{\small\ttfamily S6-\allowbreak{}LF-\allowbreak{}MAXI58} & Digi-\allowbreak{}Key (Littelfuse product) (2026). \emph{MAXI 58V Slo-\allowbreak{}Blo series 0999, part 0999080.ZXN (reseller listing)}. {\small\color{dbfSteel}https:/\allowbreak{}/\allowbreak{}www.di\allowbreak{}gikey.com (part 0999080ZXN); littelfuse.com series page returned HTTP 403; read 2026-\allowbreak{}10-\allowbreak{}29}\newline{\small 80 A at 58 V DC, 1 kA breaking capacity; series stops at 80 A (value F-\allowbreak{}FUSE-\allowbreak{}58-\allowbreak{}MAXA\allowbreak{}, reseller grade)} & {\small SPEC} \\
{\small\ttfamily S6-\allowbreak{}ESKA-\allowbreak{}58V} & ESKA (summarised from soselectronic.\allowbreak{}com listings) (2020). \emph{ESKA 343.400-\allowbreak{}58V MAXI blade fuse range}. {\small\color{dbfSteel}https:/\allowbreak{}/\allowbreak{}www.so\allowbreak{}selectronic.co\allowbreak{}m (ESKA 343.400 series); datasheet PDF HTTP 403, not read; read 2026-\allowbreak{}10-\allowbreak{}29}\newline{\small 100 A at 58 V listed 'available on request' only; stock and lead time unknown (card A-\allowbreak{}03 flip condition)} & {\small SPEC} \\
{\small\ttfamily S6-\allowbreak{}RULES-\allowbreak{}FUSE} & AIAA (2026). \emph{2026-\allowbreak{}27 DBF rules draft, 2 Sep 2026, section 3.2 propulsion: fuse type, harness and rating}. {\small\color{dbfSteel}REPO/\allowbreak{}rules/\allowbreak{}202\allowbreak{}7\_DBF\_Rules\_Dr\allowbreak{}aft.pdf (layout text analysis/\allowbreak{}data/\allowbreak{}\allowbreak{}raw/\allowbreak{}aiaa-\allowbreak{}dbf-\allowbreak{}2\allowbreak{}027-\allowbreak{}rules-\allowbreak{}draf\allowbreak{}t.layout.txt); read 2026-\allowbreak{}10-\allowbreak{}30}\newline{\small L750 one fuse per system; L754-\allowbreak{}755 blade style in a separate harness; L776-\allowbreak{}777 rating up to the pack's continuous rating, max 100 A; no voltage rating stated} & {\small RULE} \\
{\small\ttfamily S8-\allowbreak{}MILHDBK304B} & U.S. Department of Defense (data by AFPEA, Wright-\allowbreak{}Patters\allowbreak{}on AFB) (1978). \emph{MIL-\allowbreak{}HDBK-\allowbreak{}304B, Package Cushioning Design, Appendix V Graphs 10.36, 10.42, 10.48 (polyethylene 2 pcf)}. {\small\color{dbfSteel}https:/\allowbreak{}/\allowbreak{}everys\allowbreak{}pec.com/\allowbreak{}MIL-\allowbreak{}HD\allowbreak{}BK/\allowbreak{}MIL-\allowbreak{}HDBK-\allowbreak{}03\allowbreak{}00-\allowbreak{}0499/\allowbreak{}MIL\_HD\allowbreak{}BK\_304B\_1778/\allowbreak{}; read 2026-\allowbreak{}10-\allowbreak{}16 (minima, REPO analysis/\allowbreak{}data/\allowbreak{}\allowbreak{}cushion\_pe2.to\allowbreak{}ml); 2026-\allowbreak{}10-\allowbreak{}29 (curve shape, page 251 re-\allowbreak{}rendered)}\newline{\small minimum peak g per thickness at 36-\allowbreak{}48 in drops (verified); normalised curve shape g/\allowbreak{}gmin versus stress/\allowbreak{}stress\_\allowbreak{}opt read by eye from curves 3 and 4 (estimate, +-\allowbreak{}10 \%)} & {\small MEASURED} \\
{\small\ttfamily S8-\allowbreak{}RULES} & AIAA (2026). \emph{2026-\allowbreak{}27 DBF rules draft, 2 Sep 2026, sections 3.1.1-\allowbreak{}3.1.3, 3.3.4, Figure 1}. {\small\color{dbfSteel}REPO/\allowbreak{}rules/\allowbreak{}202\allowbreak{}7\_DBF\_Rules\_Dr\allowbreak{}aft.pdf (layout text analysis/\allowbreak{}data/\allowbreak{}\allowbreak{}raw/\allowbreak{}aiaa-\allowbreak{}dbf-\allowbreak{}2\allowbreak{}027-\allowbreak{}rules-\allowbreak{}draf\allowbreak{}t.layout.txt); read 2026-\allowbreak{}10-\allowbreak{}29}\newline{\small L438-\allowbreak{}453 sensor, L538-\allowbreak{}542 ring-\allowbreak{}to-\allowbreak{}bottom 15 in, L556-\allowbreak{}560 simulators, L996-\allowbreak{}1008 Ground Mission, L1002 whole inches} & {\small RULE} \\
{\small\ttfamily S8-\allowbreak{}ASTM-\allowbreak{}B777} & ASTM International; GS Alloy (2026). \emph{ASTM B777-\allowbreak{}15 (reapproved 2026), Tungsten Base High-\allowbreak{}Density Metal; GS Alloy datasheet}. {\small\color{dbfSteel}https:/\allowbreak{}/\allowbreak{}gsallo\allowbreak{}y.com/\allowbreak{}product/\allowbreak{}\allowbreak{}astm-\allowbreak{}b777-\allowbreak{}tung\allowbreak{}sten-\allowbreak{}high-\allowbreak{}dens\allowbreak{}ity-\allowbreak{}metal/\allowbreak{}; read 2026-\allowbreak{}10-\allowbreak{}29}\newline{\small density classes 1-\allowbreak{}4: 16.85-\allowbreak{}17.25, 17.15-\allowbreak{}17.85, 17.75-\allowbreak{}18.35, 18.25-\allowbreak{}18.85 g/\allowbreak{}cc; design value 17.0 g/\allowbreak{}cc} & {\small SPEC} \\
{\small\ttfamily S8-\allowbreak{}TUNGSTEN-\allowbreak{}RE\allowbreak{}TAIL} & tungsten.com (reseller) (2026). \emph{tungsten.com 1 lb bags of tungsten alloy shot and cubes, 17.6 g/\allowbreak{}cc, US\$499.99-\allowbreak{}799.\allowbreak{}99}. {\small\color{dbfSteel}https:/\allowbreak{}/\allowbreak{}shop.t\allowbreak{}ungsten.com/\allowbreak{}tu\allowbreak{}ngsten-\allowbreak{}alloy-\allowbreak{}c\allowbreak{}ubes-\allowbreak{}various-\allowbreak{}m\allowbreak{}m-\allowbreak{}sizes-\allowbreak{}indust\allowbreak{}rial-\allowbreak{}use-\allowbreak{}1-\allowbreak{}lb-\allowbreak{}\allowbreak{}bags/\allowbreak{}; read 2026-\allowbreak{}10-\allowbreak{}29}\newline{\small retail price band US\$1100-\allowbreak{}1800 per kg (computed); grade reseller. Bulk 95WNiFe rod US\$58-\allowbreak{}65/\allowbreak{}kg, 5 kg MOQ is an unverified marketplace listing (estimate)} & {\small SPEC} \\
{\small\ttfamily S8-\allowbreak{}TUNGSTEN-\allowbreak{}PU\allowbreak{}TTY} & Maximum Velocity (reseller); Pinewood Pro US\$10.95 (2026). \emph{Tungsten putty 1 oz, minimum density 10 g/\allowbreak{}cc, US\$6.25}. {\small\color{dbfSteel}https:/\allowbreak{}/\allowbreak{}maximu\allowbreak{}m-\allowbreak{}velocity.com\allowbreak{}/\allowbreak{}product/\allowbreak{}tungs\allowbreak{}ten-\allowbreak{}putty-\allowbreak{}1-\allowbreak{}ou\allowbreak{}nce; read 2026-\allowbreak{}10-\allowbreak{}29}\newline{\small putty density and price for fine-\allowbreak{}trim ballast} & {\small SPEC} \\
{\small\ttfamily S8-\allowbreak{}LEAD} & Expert Market Research (Q2 2026); Argus Media (12 Jan 2026); density from standard handbooks (not re-\allowbreak{}read) (2026). \emph{Lead ingot US\$0.89-\allowbreak{}1.17 per lb; density 11.34 g/\allowbreak{}cc}. {\small\color{dbfSteel}https:/\allowbreak{}/\allowbreak{}www.ar\allowbreak{}gusmedia.com/\allowbreak{}m\allowbreak{}etals-\allowbreak{}platform\allowbreak{}/\allowbreak{}price/\allowbreak{}assessm\allowbreak{}ent/\allowbreak{}lead-\allowbreak{}ingot\allowbreak{}-\allowbreak{}primary-\allowbreak{}99-\allowbreak{}97\allowbreak{}-\allowbreak{}-\allowbreak{}del-\allowbreak{}us-\allowbreak{}trans\allowbreak{}action-\allowbreak{}prices-\allowbreak{}\allowbreak{}PA00229370000; read 2026-\allowbreak{}10-\allowbreak{}29}\newline{\small commodity price band (reseller grade); density grade estimate} & {\small SPEC} \\
{\small\ttfamily S8-\allowbreak{}SK78-\allowbreak{}LINE} & Murray's (Robline mini spool, 880 lb); Binks Marine (Liros 410 kg) (2026). \emph{SK78 12-\allowbreak{}strand 2 mm line: break 399-\allowbreak{}410 kgf, 1.8 g/\allowbreak{}m}. {\small\color{dbfSteel}https:/\allowbreak{}/\allowbreak{}murray\allowbreak{}s.com/\allowbreak{}products\allowbreak{}/\allowbreak{}sk78-\allowbreak{}mini-\allowbreak{}spo\allowbreak{}ol-\allowbreak{}2mm-\allowbreak{}dyneema\allowbreak{}-\allowbreak{}lashing-\allowbreak{}line; read 2026-\allowbreak{}10-\allowbreak{}29}\newline{\small reference strength 3.9 kN at 2 mm, scaled with d\^{}2 for the strength member (scaling estimate)} & {\small SPEC} \\
{\small\ttfamily S8-\allowbreak{}POLOLU-\allowbreak{}4756} & Pololu; Core Electronics listing (2026). \emph{Pololu 131:1 metal gearmotor 37Dx73L mm 12 V, 64 CPR encoder (\#4756): 76 rpm, 45 kg-\allowbreak{}cm stall (extrapolated)}. {\small\color{dbfSteel}https:/\allowbreak{}/\allowbreak{}www.po\allowbreak{}lolu.com/\allowbreak{}produ\allowbreak{}ct/\allowbreak{}4756/\allowbreak{}specs; read 2026-\allowbreak{}10-\allowbreak{}29}\newline{\small added drive 37d\_131 and the 2x variant for heavier sensors (grade reseller for the search summary)} & {\small SPEC} \\
{\small\ttfamily S8-\allowbreak{}DRTRADE} & team analysis, 14 Oct 2026 (2026). \emph{DB-\allowbreak{}01 D\&R trade inputs and report (dr\_trade.toml\allowbreak{}, dr\_trade\_repor\allowbreak{}t.txt)}. {\small\color{dbfSteel}REPO/\allowbreak{}analysis/\allowbreak{}\allowbreak{}inputs/\allowbreak{}dr\_trad\allowbreak{}e.toml; read 2026-\allowbreak{}10-\allowbreak{}29}\newline{\small drives, parts masses, tension, drum sizing checks (imported read-\allowbreak{}only through s5\_tow)} & {\small COMPUTED} \\
{\small\ttfamily S8-\allowbreak{}BARROWMAN} & J. Barrowman, B. Barrowman (1967). \emph{The Theoretical Prediction of the Center of Pressure (Barrowman and Barrowman), NARAM-\allowbreak{}8}. {\small\color{dbfSteel}https:/\allowbreak{}/\allowbreak{}www.ap\allowbreak{}ogeerockets.co\allowbreak{}m/\allowbreak{}education/\allowbreak{}do\allowbreak{}wnloads/\allowbreak{}Barrow\allowbreak{}man\_Report.pdf; read formulae applied from the method; paper not re-\allowbreak{}read this session}\newline{\small nose CNa 2, CP 0.466 L (ogive); 4-\allowbreak{}fin trapezoid CNa and CP; subsonic} & {\small METHOD} \\
{\small\ttfamily S8-\allowbreak{}NACA-\allowbreak{}TM832} & NACA (1937). \emph{NACA TM 832 towed-\allowbreak{}body guidance (heavy, compact, damped, small fins far aft), as cited in brief section 5}. {\small\color{dbfSteel}not re-\allowbreak{}read; cited in the orchestrator brief; read not read by S8}\newline{\small design guidance only; no number taken from it} & {\small METHOD} \\
\end{dbflongtable}

